\begin{pdfdiagram}[x=1cm,y=1cm]
  % 0--0.9 V 的完整绘图区。
  \fill[BookCodeBg] (0,0) rectangle (5.8,5.8);
  \draw[BookRule, line width=0.45pt] (0,0) rectangle (5.8,5.8);
  \draw[BookRule, line width=0.6pt, ->] (0,0) -- (6.25,0)
    node[right, hw label] {$V_Q$ (V)};
  \draw[BookRule, line width=0.6pt, ->] (0,0) -- (0,6.25)
    node[above, hw label] {$V_{\overline{Q}}$ (V)};
  \foreach \p/\t in {0/0,2.9/0.45,5.8/0.9} {
    \draw[BookMuted, line width=0.45pt] (\p,0.07) -- (\p,-0.07);
    \draw[BookMuted, line width=0.45pt] (0.07,\p) -- (-0.07,\p);
    \node[hw label, anchor=north] at (\p,-0.10) {\t};
    \node[hw label, anchor=east] at (-0.12,\p) {\t};
  }
  \draw[BookMuted, dashed, line width=0.45pt] (0,0) -- (5.8,5.8);
  \node[hw label, anchor=south east] at (5.65,5.48) {镜像轴 $V_Q=V_{\overline Q}$};
  \begin{scope}
    \clip (0,0) rectangle (5.8,5.8);
    \draw[BookRed, line width=1.0pt]
      (0.35,5.25) .. controls (1.20,5.25) and (1.75,4.85) .. (2.55,2.90)
      .. controls (3.28,1.00) and (3.85,0.45) .. (5.35,0.45);
    \draw[BookAccent, line width=1.0pt]
      (5.25,5.35) .. controls (3.85,5.28) and (3.20,4.75) .. (2.55,2.90)
      .. controls (1.78,1.00) and (1.18,0.45) .. (0.35,0.35);
  \end{scope}
  \draw[BookAmber, line width=0.9pt, fill=BookAmber!8] (0.82,3.18) rectangle (2.22,4.58);
  \draw[BookAmber, line width=0.9pt, fill=BookAmber!8] (3.18,1.02) rectangle (4.58,2.42);
  \node[BookAmber, font=\small\bfseries] at (1.52,3.88) {SNM$_1$};
  \node[BookAmber, font=\small\bfseries] at (3.88,1.72) {SNM$_2$};
  \fill[BookRed] (0.52,5.05) circle (0.075);
  \fill[BookRed] (5.05,0.52) circle (0.075);
  \fill[BookMuted] (2.55,2.90) circle (0.065);
  \node[BookRed, font=\tiny\bfseries, anchor=west] at (0.65,5.25) {稳态 1};
  \node[BookRed, font=\tiny\bfseries, anchor=east] at (4.95,0.82) {稳态 2};
  \node[hw label, anchor=west] at (2.72,2.72) {亚稳态};
  \draw[BookRed, line width=1pt] (6.65,5.25) -- (7.35,5.25);
  \node[hw label, anchor=west] at (7.48,5.25) {反相器 1 的 VTC};
  \draw[BookAccent, line width=1pt] (6.65,4.72) -- (7.35,4.72);
  \node[hw label, anchor=west] at (7.48,4.72) {反相器 2 的镜像 VTC};
  \draw[BookAmber, line width=0.9pt] (6.65,4.05) rectangle (7.35,4.55);
  \node[hw label, anchor=west, text width=3.2cm] at (7.48,4.30)
    {每个“眼”中能放入的最大正方形};
  \node[hw label, anchor=west, text width=4.1cm] at (6.65,3.20)
    {读图步骤：\\1. 两条 VTC 围出两个眼；\\2. 分别找最大内切正方形；\\3. 取较小边长作为单元 SNM。};
  \node[hw label, anchor=west, text=BookAmber] at (6.65,1.50)
    {$\mathrm{SNM}=\min(\mathrm{SNM}_1,\mathrm{SNM}_2)$};
\end{pdfdiagram}
